\documentclass[11pt]{article}
\usepackage{fontspec}
\usepackage[french]{babel}
\usepackage{amsmath}
\usepackage{amssymb}
\usepackage{amsthm}
\usepackage{mathrsfs}
\usepackage{tikz-cd}
\usepackage{enumitem}

\newcommand{\printedpage}[1]{\marginpar{\footnotesize p.~#1}}
\newcommand{\msp}[1]{\marginpar{\footnotesize [#1]}}
\newcommand{\circled}[1]{\tikz[baseline=(c.base)]\node[draw,circle,inner sep=1pt](c){#1};}
\newcommand{\Cat}{\mathrm{Cat}}
\newcommand{\Ob}{\mathrm{Ob}\,}
\newcommand{\hot}{\mathrm{hot}}

\begin{document}

% ===================== page 1 (PDF) =====================
\printedpage{[non numérotée]}
% head: (none printed)

\begin{center}
{\sffamily\LARGE LES DÉRIVATEURS}

\bigskip

{\large\textbf{ALEXANDRE GROTHENDIECK}}

\vspace{3em}

{\large\textbf{Chapitre XVI}}

\medskip

{\large\textbf{Les localiseurs fondamentaux}}
\end{center}

\clearpage

% ===================== page 2 (PDF) =====================
\printedpage{[non numérotée]}
% head: (none printed)

Ce texte a été déchiffré et transcrit en \LaTeXe{} par M. Künzer. Il a été édité par M. Künzer, J. Malgoire, et G. Maltsiniotis. La transcription est aussi fidèle que possible au manuscrit. Seules quelques corrections mineures évidentes ont été effectuées. Pour les rares ajouts ou commentaires des éditeurs, ainsi que pour la numérotation originale des pages du manuscrit, les caractères de machine à écrire \texttt{[typewriter]} entre crochets sont utilisés.

Cette édition est provisoire. Les remarques, commentaires et corrections sont bienvenus. Envoyer un message à :

\textbf{georges.maltsiniotis@imj-prg.fr}

\hfill G. Maltsiniotis\hspace*{2cm}

\clearpage

% ===================== page 3 (PDF) =====================
\printedpage{1}
% head: A. Grothendieck : Dérivateurs, Ch. XVI, Édition des Universités de Montpellier II et Paris VII
% foot: Version 18 mars 2019

\msp{page 1}

\section*{1\quad Axiomes Loc(1) à Loc(3) (et variantes) : localiseurs fondamentaux}

J’ai envie de reprendre l’exposé des axiomes utiles pour un localiseur fondamental
\[
W \subset \mathrm{Fl}(\Cat),
\]
en me guidant cette fois plus sur la signification géométrique des axiomes, que sur les alias $\pm$ techniques qu’on rencontre (ou ne rencontre pas) pour les vérifier. Je révise donc dans cet esprit la présentation de XV §9.

\begin{description}
\item[Loc (1)] (\textit{saturation}).
\begin{itemize}
\item[\textit{a})] Iso $\subset W$.
\item[\textit{b})] Si deux parmi $f$, $g$, $gf$ sont dans $W$, aussi le troisième.
\item[\textit{c})] Si $gf$ et $fg$ sont dans $W$, $f$ et $g$ aussi.
\end{itemize}
\end{description}

Les conditions (\textit{a}) et (\textit{b}) toutes seules, je vais les appeler (pour un $W \subset \mathrm{Fl}\,\mathcal{M}$ général) \textit{saturation faible} du localiseur. Les conditions (\textit{a}), (\textit{b}), (\textit{c}), \textit{saturation modérée} (“mild saturation”). Autres variantes :

\begin{description}
\item[Loc (1 bis)] Les conditions (\textit{a}), (\textit{b}), (\textit{c}) du haut, plus
\begin{itemize}
\item[\textit{d})] $W$ stable par facteurs directs.
\end{itemize}
\end{description}

C’est bien sûr impliqué par :

\begin{description}
\item[Loc (1 ter)] $W$ est \textit{fortement saturé, i.e.} si $f \in \Cat$ est telle que $\hot_W(f)$ est un isomorphisme, alors $f \in W$.
\item[Loc (2)] (\textit{Objet final}). Si $X$ dans Cat a un objet final, alors $X$ est $W$-asphérique, \textit{i.e.} $X \longrightarrow e$ est dans $W$.
\end{description}

Parfois, cet axiome est conséquence de la variante affaiblie suivante.

\begin{description}
\item[Loc (2$_0$)] (Axiome du segment). $\boldsymbol{\Delta}^1$ est $W$-asphérique.
\end{description}

\msp{page 2}

\begin{description}
\item[Loc (3)] (\textit{localisation,} forme faible). Soit $f : X \longrightarrow S$ une flèche dans Cat $W$\textit{-asphérique, i.e.} telle que $\forall\ s \in \Ob S$, la catégorie localisée $X/s$ sur $S$ en $s$, $X/s = X \times_S S/s$, soit $W$-asphérique. Alors $f \in W$.\footnote[1]{Dans une version finale, il faudrait intervertir les notations pour Loc (3) et pour Loc (3 bis), car ce dernier m’apparaît comme le plus important des deux - il est satisfait pour tous les localiseurs qui ont été utilisées (à ma connaissance), et est valable pour tout localiseur de la forme $W_{\mathbf{D}}$, où $\mathbf{D}$ est un dérivateur (ou seulement un prédérivateur satisfaisant Der (4,5) \texttt{[il faut aussi Der (2) ; la numérotation des axiomes est celle du chapitre I, p. 68]}).} % mark printed as superscript (1); [sic] utilisées -> utilisés, agreement with "localiseurs"
\end{description}

\textbf{Rappel 1.} Loc (1, 2, 3) impliquent que l’on a
\[
W = W^\circ \qquad \textit{i.e. } f \in W \iff f^\circ \in W,
\]
\printedpage{2}
% head: A. Grothendieck : Dérivateurs, Ch. XVI, Édition des Universités de Montpellier II et Paris VII
% foot: Version 18 mars 2019
où on dénote par $X \longmapsto X^\circ$, $f \longmapsto f^\circ$ le foncteur de passage à la catégorie opposée, foncteur
\[
X \longmapsto X^\circ : \Cat \xrightarrow{\ \sim\ } \Cat
\]
qui est un \textit{automorphisme} de Cat (donc commute à toutes $\varinjlim$ et $\varprojlim$ ). Il s’ensuit que l’ensemble d’axiomes Loc (1, 2, 3) équivaut à Loc (1, 2$^{\mathrm{opp}}$, 3$^{\mathrm{opp}}$), où Loc (2$^{\mathrm{opp}}$) est l’axiome “dual” de Loc (2), savoir que tout objet de Cat ayant un objet initial est $W$-asphérique, et de même Loc (3$^{\mathrm{opp}}$) l’axiome que toute $f : X \longrightarrow S$ $W$\textit{-coasphérique, i.e.} telle que les catégories colocalisées $s\backslash X$ de $X$ sur $S$ soient $W$-asphériques, est dans $W$.

\textbf{Définition 1.1.} On appelle \textit{localiseur fondamental} une partie $W$ de Fl(Cat) satisfaisant les axiomes Loc (1, 2, 3) (\textit{i.e.} saturation modérée, axiome de l’objet final, axiome de localisation).

\msp{page 3}

\textbf{Rappel 1.2.} Si $W$ est un localiseur fondamental, alors tout \textit{homotopisme} $f : X \longrightarrow Y$ de Cat est dans $W$ (c’est pour cela essentiellement que l’on a imposé (\textit{c}) dans la saturation modérée Loc (1)). Si $f, g : X \rightrightarrows Y$ sont homotopes (\textit{i.e.} appartiennent à une même composante connexe de $\underline{\mathrm{Hom}}(X, Y)$), alors $\hot_W(f) = \hot_W(g)$ et $f \in W \iff g \in W$.

Voici la forme forte de Loc (3).

\begin{description}
\item[Loc (3 bis)] (\textit{Localisation,} forme forte). Soit $S$ dans Cat et $f : X \longrightarrow Y$ un morphisme dans $\Cat/S$ tel que “$f$ soit une $W$-équivalence localement sur $S$”, \textit{i.e.} pour tout $s \in \Ob S$, le morphisme induit
\[
f/s : X/s \longrightarrow Y/s
\]
pour les catégories localisées de $X$, $Y$ sur $S$ en $s$, soit dans $W$. Alors $f \in W$.
\end{description}

L’axiome Loc (3) n’est autre que le cas particulier de Loc (3 bis) relatif au cas où $Y = S$.

\textbf{Proposition 1.3.} \textit{Soit $W \subset \mathrm{Fl}(\mathcal{M})$. Alors on a les équivalences suivantes}
\[
\begin{aligned}
\text{Loc (1, 2, 3)} &\iff \text{Loc (1, $2_0$, $3'$)}\\
\text{Loc (1, 2, 3 bis)} &\iff \text{Loc (1, $2_0$, $3'$ bis),}
\end{aligned}
\]
\textit{où les axiomes} Loc ($3'$) \textit{et} Loc ($3'$ bis) \textit{se formulent ainsi :}

\msp{page 4}

\begin{description}
\item[Loc (3$'$)] \circled{\textit{a}} (Axiome des Cat-fibrations, forme directe faible). Soient $S$ dans Cat, $X$ une catégorie fibrée scindée sur $S$, telle que les fibres $X_s$ ($s \in \Ob S$) soient $W$-asphériques. Alors $X \longrightarrow S$ est dans $W$.

\circled{\textit{b}} Soit $X$ dans $\Cat/S$, alors $i_X : X \longrightarrow \Phi_S(X)$ est dans $W$\footnote[2]{(\textit{b}) concerne la \textit{comparaison} d’une $X$ quelconque sur $S$, avec la catégorie fibrée scindée $\Phi_S(X)$ sur $S$.}. % mark printed as superscript (2)
\end{description}

\clearpage

% ===================== page 5 (PDF) =====================
\printedpage{3}
% head: A. Grothendieck : Dérivateurs, Ch. XVI, Édition des Universités de Montpellier II et Paris VII
% foot: Version 18 mars 2019

\begin{description}
\item[Loc (3$'$ bis)] \circled{\textit{a}} (Axiome des Cat-fibrations, forme directe forte). Soient $S$ dans Cat, $X$, $Y$ deux catégories fibrées scindées sur $S$, $f : X \longrightarrow Y$ un $S$-foncteur cartésien compatible aux scindages, tel que $\forall\ s \in \Ob S$, le foncteur induit sur les fibres
\[
f_s : X_s \longrightarrow Y_s
\]
soit dans $W$. Alors $f \in W$.

\circled{\textit{b}} Comme Loc ($3'$) (\textit{b}).
\end{description}

\textbf{N.B.} En fait, moyennant Loc (1, 2), les énoncés Loc ($3'$), Loc ($3'$ bis) impliquent les énoncés similaires sans scindages, et sans même supposer $f$ cartésien dans Loc ($3'$ bis), \textit{cf.} démonstration ci-dessous.

\textsc{Démonstration.} \textit{a}) Loc (1, 2, 3) $\Longrightarrow$ Loc (1, $2_0$, $3'$). Il suffit de prouver Loc ($3'$). Prouvons Loc ($3'$ $a$). Pour ceci, sous les conditions de Loc ($3'$ $a$), prouvons que $X \longrightarrow S$ est $W$-coasphérique, \textit{i.e.} que les catégories
\[
s\backslash X = X \times_S s\backslash S
\]
sont $W$-asphériques. Il en résultera bien que $f \in W$ en vertu de Loc (3) appliqué à $f^\circ$ et de $W = W^\circ$. Mais
\msp{page 5}
considérons l’inclusion canonique
\[
i_s : X_s \lhook\joinrel\longrightarrow s\backslash X,
\]
il résulte du fait que $p : X \longrightarrow S$ est Cat-fibrant que ce foncteur est $W$-asphérique (il a même un adjoint à droite, \textit{i.e.} les catégories $X_s/u$, où $u \in \Ob(s\backslash X)$, ont un objet final). Donc par Loc (3) on a $i_s \in W$, et comme $X_s$ est $W$-asphérique par hypothèse, de même $s\backslash X$, q.e.d.

Prouvons Loc ($3'$ $b$) : On note que $i_X : X \longrightarrow \Phi_S(X)$ est $W$\textit{-coasphérique} (il a même un adjoint à gauche), donc $i_X^\circ$ est $W$-asphérique, donc dans $W$, donc $i_X^\circ \in W^\circ$ car $W = W^\circ$. % [sic] i_X^\circ \in W^\circ -> i_X \in W (or i_X \in W^\circ), apparently intended: from i_X^\circ \in W one gets i_X \in W^\circ = W

$a$ bis) Loc (1, 2, 3 bis) $\Longrightarrow$ Loc (1, $2_0$, $3'$ bis). Cette fois il faut prouver seulement Loc ($3'$ bis). Pour (\textit{b}), c’est déjà fait. Reste à prouver (\textit{a}). Si $f$ est comme dans l’énoncé Loc ($3'$ bis $a$), prouvons que pour tout $s$ dans $S$, le morphisme induit
\[
s\backslash f : s\backslash X \longrightarrow s\backslash Y
\]
est dans $W$. Comme $W = W^\circ$, on aura donc $(s\backslash f)^\circ \in W$, et appliquant Loc (3 bis) à $f^\circ$, on trouve $f^\circ \in W$, donc $f \in W$, et on aura fini. Mais $s\backslash f$ s’insère dans le carré commutatif
\msp{page 6}
\[
\begin{tikzcd}
X_s \arrow[r, "f_s"] \arrow[d, "i_s^X"'] & Y_s \arrow[d, "i_s^Y"] \\
s\backslash X \arrow[r, "s\backslash f"] & s\backslash Y,
\end{tikzcd}
\]

\end{document}
